% -----------------------------------------------------------------------------------------------------
\section{区域分析}\label{sec:sudakov-regions}

\cref{prop:sudakov-oneloop}的双极点从何而来？\cref{ch:ir}的区域方法（\cref{def:ir-regions}）回答这个问题：按圈动量的标度划分区域，在每个区域内 Taylor 展开、
在整个积分域上用维数正规化积分、丢掉无标度积分，再求和。三个一圈例子说明了它在 Sudakov 问题中的全部要点。我们研究标量三角形
\begin{equation}\label{eq:sudakov-triangle}
I=\ee^{\varepsilon\gamma_E}\mu^{2\varepsilon}\int\!\frac{\dd^Dk}{\ii\pi^{D/2}}\,
\frac{1}{(-k^2-\ii0)\,(-(k+p_1)^2-\ii0)\,(-(k+p_2)^2-\ii0)},
\end{equation}
领头对数精度下 QCD 顶点修正正比于 $p_1\cdot p_2\,I$（\cref{ex:sudakov-ll}）。

\begin{example}[在壳无质量形状因子 = 硬区]\label{ex:sudakov-onshell-hard}
取 $p_1^2=p_2^2=0$，$(p_1-p_2)^2=-Q^2$。由 Feynman 参数表示 $I=\ee^{\varepsilon\gamma_E}\mu^{2\varepsilon}\Gamma(1+\varepsilon)\int\!\dd F\,(Q^2xy)^{-1-\varepsilon}$ 与 Dirichlet 积分，
\begin{equation}\label{eq:sudakov-onshell-H}
Q^2I=\Big(\frac{\mu^2}{Q^2}\Big)^{\varepsilon}\frac{\ee^{\varepsilon\gamma_E}\Gamma(1+\varepsilon)\Gamma^2(-\varepsilon)}{\Gamma(1-2\varepsilon)}
\overset{\mu=Q}{=}\frac1{\varepsilon^2}-\frac{\pi^2}{12}-\frac73\zeta_3\varepsilon+O(\varepsilon^2)
% VF-REG-10 (H)
\end{equation}
（在 $\Real\varepsilon<0$ 收敛：极点是红外的）。被积函数中没有小量，硬区 $k\sim Q(1,1,1)$ 的展开就是原被积函数；其余区域都无标度，例如与 $p_1$ 共线的区域把
$-(k+p_2)^2$ 换成 $-2p_2^-k^+$ 后，被积函数在 $k^-\to\kappa^2k^-$、$\mathbf k_T\to\kappa\mathbf k_T$ 下齐次。所以在壳积分的红外极点是\emph{无标度低能区域的紫外极点}：
低能区域 $\propto\frac1{\varepsilon_{\rm UV}}-\frac1{\varepsilon_{\rm IR}}=0$，其紫外部分与硬区的红外极点相消。这一机制适用于整个顶点，也是\cref{thm:sudakov-ir-poles}中 $\mathbf Z$ 因子的起源。
\uses{prop:sudakov-oneloop,def:ir-regions}
\end{example}

\begin{example}[离壳 Sudakov 三角形的区域展开]\label{ex:sudakov-offshell}
取 $p_1^2=p_2^2=-P^2<0$，$(p_1-p_2)^2=-Q^2$，$t\equiv P^2/Q^2\ll1$，$\mathcal L\equiv\ln(Q^2/P^2)$。此时 $\mathcal F=P^2\alpha_0(\alpha_1+\alpha_2)+Q^2\alpha_1\alpha_2>0$，
$I$ 在 $D=4$ 中有限且为实数；令 $\alpha_0=1-r$、$\alpha_1=ru$、$\alpha_2=r(1-u)$ 并对 $r$ 积分，
\begin{equation}\label{eq:sudakov-offshell-exact}
Q^2I\big|_{\varepsilon=0}=\int_0^1\!\dd u\,\frac{\ln\big(u(1-u)/t\big)}{u(1-u)-t}
=\frac{4\,\mathrm{Li}_2(-\rho t)+\ln^2(\rho t)+\pi^2/3}{\sqrt{1-4t}},\qquad\rho=\frac{2}{1-2t+\sqrt{1-4t}},
% VF-REF-2
\end{equation}
即 Usyukina--Davydychev 函数 $\Phi^{(1)}(t,t)$\cite{UsyukinaDavydychev1993a}；其展开为
\begin{equation}\label{eq:sudakov-offshell-nlp}
Q^2I=\mathcal L^2+\frac{\pi^2}3+t\Big(2\mathcal L^2-4\mathcal L-4+\frac{2\pi^2}3\Big)+t^2\big(6\mathcal L^2-14\mathcal L-11+2\pi^2\big)+O(t^3\mathcal L^2).
% VF-REF-2 (expansion; binding sign convention L = ln(Q^2/m^2))
\end{equation}
$Q^2\gg P^2$ 不能用欧氏外动量实现（欧氏几何中 $|p_1-p_2|^2\le4P^2$），实动量是类空 Minkowski 矢量 $p_1=(p_1^+,-\frac{P^2}{2p_1^+},\mathbf 0_T)$、$p_2=(-\frac{P^2}{2p_2^-},p_2^-,\mathbf 0_T)$。

\emph{区域。}$\lambda^2\sim t$ 时有四个贡献区域（\cref{fig:sudakov-ff}；标度按 $(k^+,k^-,k_T)$）：硬区 $Q(1,1,1)$、共线区 $C_1$：$Q(1,\lambda^2,\lambda)$、$C_2$：$Q(\lambda^2,1,\lambda)$、
超软区 $S$：$Q(\lambda^2,\lambda^2,\lambda^2)$；记 $\mathcal N(\varepsilon)=\ee^{\varepsilon\gamma_E}\Gamma(1+\varepsilon)\Gamma^2(-\varepsilon)/\Gamma(1-2\varepsilon)$，
\begin{equation}\label{eq:sudakov-regions}
Q^2H=\Big(\frac{\mu^2}{Q^2}\Big)^{\varepsilon}\mathcal N,\qquad Q^2C_{1,2}=-\Big(\frac{\mu^2}{P^2}\Big)^{\varepsilon}\mathcal N,\qquad
Q^2S=-\Big(\frac{\mu^2Q^2}{P^4}\Big)^{\varepsilon}\ee^{\varepsilon\gamma_E}\Gamma(1+\varepsilon)\Gamma(\varepsilon)\Gamma(-\varepsilon).
% VF-REG-10
\end{equation}
以 $C_1$ 为例：在 $-(k+p_2)^2$ 中只保留 $-2p_2^-k^+$，用参数 $x$ 合并 $-k^2$ 与 $-(k+p_1)^2$（得 $-\ell^2+x(1-x)P^2$，$\ell=k+xp_1$），以 $\frac{1}{A^2B}=2\int_0^\infty\frac{\dd\eta}{(A+\eta B)^3}$
并入线性分母并作类光平移 $\ell\to\ell+\eta p_2^-n_B$，得 $C_1=\ee^{\varepsilon\gamma_E}\Gamma(1+\varepsilon)\mu^{2\varepsilon}\int_0^1\!\dd x\int_0^\infty\!\dd\eta\,[x(1-x)P^2+\eta xQ^2]^{-1-\varepsilon}$
（$\eta$ 积分在 $\Real\varepsilon>0$、$x$ 积分在 $\Real\varepsilon<0$ 收敛，各自解析延拓）。超软区用两个 Schwinger 参数，由
$-k^2-2\eta_1p_1^+k^--2\eta_2p_2^-k^+=-(k+\eta_1p_1^+n_A+\eta_2p_2^-n_B)^2+\eta_1\eta_2Q^2$ 化为 $\iint\dd\eta_1\dd\eta_2[\eta_1\eta_2Q^2+(\eta_1+\eta_2)P^2]^{-1-\varepsilon}$，令 $\eta_i=ts_i$ 逐次积分即得。
软区 $Q(\lambda,\lambda,\lambda)$ 的分母齐次，无标度；Glauber 区 $Q(\lambda^2,\lambda^2,\lambda)$ 中 $k^\pm$ 各只出现在一个分母里，加解析正规子后为零。
两两重叠（如 $T_ST_{C_1}\propto\int_0^\infty\dd\eta_1\,\eta_1^{-1-\varepsilon}$）无标度：维数正规化中\emph{不需要零仓扣除}。

\emph{求和。}取 $\mu=Q$，$Q^2H=\frac1{\varepsilon^2}-\frac{\pi^2}{12}$，$Q^2C_{1,2}=-\frac1{\varepsilon^2}-\frac{\mathcal L}\varepsilon-\frac{\mathcal L^2}2+\frac{\pi^2}{12}$，
$Q^2S=\frac1{\varepsilon^2}+\frac{2\mathcal L}\varepsilon+2\mathcal L^2+\frac{\pi^2}4$（均差 $O(\varepsilon)$），
\begin{equation}\label{eq:sudakov-regions-sum}
Q^2(H+C_1+C_2+S)=\mathcal L^2+\frac{\pi^2}3+\varepsilon\Big(\mathcal L^3+\frac{2\pi^2}3\mathcal L+2\zeta_3\Big)+O(\varepsilon^2).
% VF-REG-10
\end{equation}
$1/\varepsilon^2$ 的系数 $1-2+1=0$；对一般的 $\mu$，$1/\varepsilon$ 的系数 $\ln\frac{\mu^2}{Q^2}-2\ln\frac{\mu^2}{P^2}+\ln\frac{\mu^2Q^2}{P^4}=0$。$H$ 的极点是红外的，$S$ 的是紫外的，$C_i$ 兼有：
硬区的红外极点被低能区的紫外极点抵消。常数 $-\frac{\pi^2}{12}+\frac{\pi^2}6+\frac{\pi^2}4=\frac{\pi^2}3$，与式~\eqref{eq:sudakov-offshell-nlp} 的领头幂次相同；
各区域展开到 $t$ 与 $t^2$ 阶时，次领头与次次领头系数也逐项重现，且每一阶的极点分别相消。不等虚度时领头幂次为 $\ln\frac{Q^2}{P_1^2}\ln\frac{Q^2}{P_2^2}+\frac{\pi^2}3$。% check: s2_regions.py (C); s8_mb_regions.py (Mellin-Barnes pole families n=0,1,2)
硬区、共线区与超软区的结果与\cite[第~2.2~节，式~(2.39)--(2.44)]{BecherBroggioFerroglia2015}去掉 $\ee^{\varepsilon\gamma_E}$ 后逐一相同（那里另一共线区由对称性得到）；方法源于\cite{BenekeSmirnov1998}（针对阈展开）。
Passarino--Veltman 函数（Minkowski 传播子，测度 $\dd^Dk/(\ii\pi^{D/2})$）满足 $q^2C_0=Q^2I$。% VF-REG-14
这个四侧面的 Minkowski 例子不在区域方法的任何一般证明之内（\cref{prop:ir-regions-minkowski}，[P]）。同一三角形的 Collins 减除项与精确余项（余项在次领头幂次没有对数）
见第~9~章\cref{ex:method-sudakov-subtraction}与附录~N 的\cref{prop:appN-triangle-exact}，积分范数下的验证见\cref{prop:appN-triangle-cs}，两圈的九个区域见\cref{fig:method-hasse}。
\uses{def:ir-regions,prop:ir-regions-minkowski}
\end{example}

%%% BEGIN fig:sudakov-ff
\begin{figure}[htbp]
\centering
\begin{tikzpicture}[fd,
  t2 hc/.style={draw=black,line width=1.15pt},
  t2 c1/.style={draw=blue!60!black},
  t2 c2/.style={draw=orange!80!black},
  t2 us/.style={draw=green!45!black},
  t2 v/.style={font=\scriptsize,inner sep=1pt},
  t2 lab/.style={font=\small}]
  % ======================= master graph with momentum routing =======================
  \begin{scope}
    \coordinate (V) at (0,0); \coordinate (E1) at (2.45,1.35); \coordinate (E2) at (2.45,-1.35);
    \coordinate (G1) at ($(V)!0.58!(E1)$); \coordinate (G2) at ($(V)!0.58!(E2)$);
    \draw[photon] (-1.25,0) -- (V);
    \draw[quark] (V) -- (G1); \draw[quark] (G1) -- (E1);
    \draw[quark] (E2) -- (G2); \draw[quark] (G2) -- (V);
    \draw[gluon] (G2) -- (G1);
    \foreach \p in {V,G1,G2} \node[vertex] at (\p) {};
    \path[mom={$q$},mom sep=0.28cm] (-1.25,0) -- (V);
    \path[mom={$p_1-k$},mom pos=0.5] (V) -- (G1);
    \path[mom'={$p_2+k$},mom pos=0.5] (V) -- (G2);
    \path[mom'={$k$},mom sep=0.3cm] (G2) -- (G1);
    \path[mom={$p_1$},mom pos=0.55] (G1) -- (E1);
    \path[mom'={$p_2$},mom pos=0.55] (G2) -- (E2);
    \node[font=\footnotesize,align=center] at (0.75,-2.05) {$p_1^2\sim p_2^2\sim\lambda^2Q^2\ll Q^2=q^2$};
  \end{scope}
  % region panels: smaller copies (scale 0.74) of the same graph, internal lines coloured by their virtuality class
  % #1 = shift, line styles for (V-G1), (V-G2), gluon; labels
  % ---------------- hard ----------------
  \begin{scope}[xshift=4.85cm,yshift=1.05cm,scale=0.74,every node/.append style={transform shape=false}]
    \coordinate (V) at (0,0); \coordinate (E1) at (2.45,1.35); \coordinate (E2) at (2.45,-1.35);
    \coordinate (G1) at ($(V)!0.58!(E1)$); \coordinate (G2) at ($(V)!0.58!(E2)$);
    \fill[gray!28] (V) -- (G1) -- (G2) -- cycle;
    \draw[photon] (-1.2,0) -- (V);
    \draw[quark,t2 hc] (V) -- (G1); \draw[quark] (G1) -- (E1);
    \draw[quark] (E2) -- (G2); \draw[quark,t2 hc] (G2) -- (V);
    \draw[gluon,line width=1.0pt] (G2) -- (G1);
    \node[t2 v,anchor=south east] at ($(V)!0.5!(G1)$) {$Q^2$};
    \node[t2 v,anchor=north east] at ($(V)!0.5!(G2)$) {$Q^2$};
    \node[t2 v,anchor=west] at ($(G1)!0.5!(G2)+(0.2,0)$) {$Q^2$};
    \node[font=\footnotesize,anchor=north] at (0.75,-1.45) {硬区：$k\sim Q(1,1,1)$};
  \end{scope}
  % ---------------- collinear to p1 ----------------
  \begin{scope}[xshift=9.35cm,yshift=1.05cm,scale=0.74,every node/.append style={transform shape=false}]
    \coordinate (V) at (0,0); \coordinate (E1) at (2.45,1.35); \coordinate (E2) at (2.45,-1.35);
    \coordinate (G1) at ($(V)!0.58!(E1)$); \coordinate (G2) at ($(V)!0.58!(E2)$);
    \draw[photon] (-1.2,0) -- (V);
    \draw[quark,t2 c1] (V) -- (G1); \draw[quark] (G1) -- (E1);
    \draw[quark] (E2) -- (G2); \draw[quark,t2 hc] (G2) -- (V);
    \draw[gluon,t2 c1] (G2) -- (G1);
    \node[t2 v,anchor=south east,text=blue!50!black] at ($(V)!0.5!(G1)$) {$\lambda^2Q^2$};
    \node[t2 v,anchor=north east] at ($(V)!0.5!(G2)$) {$Q^2$};
    \node[t2 v,anchor=west,text=blue!50!black] at ($(G1)!0.5!(G2)+(0.2,0)$) {$\lambda^2Q^2$};
    \node[font=\footnotesize,anchor=north] at (0.75,-1.45) {$1$ 共线：$k\sim Q(1,\lambda^2,\lambda)$};
  \end{scope}
  % ---------------- collinear to p2 ----------------
  \begin{scope}[xshift=4.85cm,yshift=-1.95cm,scale=0.74,every node/.append style={transform shape=false}]
    \coordinate (V) at (0,0); \coordinate (E1) at (2.45,1.35); \coordinate (E2) at (2.45,-1.35);
    \coordinate (G1) at ($(V)!0.58!(E1)$); \coordinate (G2) at ($(V)!0.58!(E2)$);
    \draw[photon] (-1.2,0) -- (V);
    \draw[quark,t2 hc] (V) -- (G1); \draw[quark] (G1) -- (E1);
    \draw[quark] (E2) -- (G2); \draw[quark,t2 c2] (G2) -- (V);
    \draw[gluon,t2 c2] (G2) -- (G1);
    \node[t2 v,anchor=south east] at ($(V)!0.5!(G1)$) {$Q^2$};
    \node[t2 v,anchor=north east,text=orange!60!black] at ($(V)!0.5!(G2)$) {$\lambda^2Q^2$};
    \node[t2 v,anchor=west,text=orange!60!black] at ($(G1)!0.5!(G2)+(0.2,0)$) {$\lambda^2Q^2$};
    \node[font=\footnotesize,anchor=north] at (0.75,-1.45) {$2$ 共线：$k\sim Q(\lambda^2,1,\lambda)$};
  \end{scope}
  % ---------------- ultrasoft ----------------
  \begin{scope}[xshift=9.35cm,yshift=-1.95cm,scale=0.74,every node/.append style={transform shape=false}]
    \coordinate (V) at (0,0); \coordinate (E1) at (2.45,1.35); \coordinate (E2) at (2.45,-1.35);
    \coordinate (G1) at ($(V)!0.58!(E1)$); \coordinate (G2) at ($(V)!0.58!(E2)$);
    \draw[photon] (-1.2,0) -- (V);
    \draw[quark,t2 c1] (V) -- (G1); \draw[quark] (G1) -- (E1);
    \draw[quark] (E2) -- (G2); \draw[quark,t2 c2] (G2) -- (V);
    \draw[gluon,t2 us] (G2) -- (G1);
    \node[t2 v,anchor=south east,text=blue!50!black] at ($(V)!0.5!(G1)$) {$\lambda^2Q^2$};
    \node[t2 v,anchor=north east,text=orange!60!black] at ($(V)!0.5!(G2)$) {$\lambda^2Q^2$};
    \node[t2 v,anchor=west,text=green!35!black] at ($(G1)!0.5!(G2)+(0.2,0)$) {$\lambda^4Q^2$};
    \node[font=\footnotesize,anchor=north] at (0.75,-1.45) {超软：$k\sim Q(\lambda^2,\lambda^2,\lambda^2)$};
  \end{scope}
\end{tikzpicture}
\caption{单圈离壳 Sudakov 形状因子的按区域展开（此处画成类时情形 $\gamma^*(q)\to q(p_1)\,\bar q(p_2)$；$p_1$ 沿 $n$（$+$ 方向），$p_2$ 沿 $\bar n$）。
左：动量走向，外线离壳 $p_1\sim Q(1,\lambda^2,\lambda)$、$p_2\sim Q(\lambda^2,1,\lambda)$，标度均按 $(k^+,k^-,k_T)$ 书写。
右：圈动量 $k$ 的四个贡献区域，内线旁标出其虚度的量级：硬区（三条内线均离壳 $\sim Q^2$，整个圈收缩为一点，灰色）、
$1$ 共线区（$k$ 与 $p_1-k$ 构成沿 $p_1$ 的喷注，蓝色；$p_2+k$ 为硬线）、$2$ 共线区（橙色，镜像）、超软区（胶子虚度 $\sim\lambda^4Q^2$，绿色；两条夸克线虚度 $\sim\lambda^2Q^2$）。
软区 $k\sim Q(\lambda,\lambda,\lambda)$ 给出无标度积分，在维数正规化中为零。各区域单独含有 $1/\varepsilon^2$、$1/\varepsilon$ 极点：对标量三角形积分，这些极点在求和中完全相消，给出有限的原积分；
对完整的 QCD 顶点，求和后只剩原积分本身的紫外发散，由重整化消去。}
\label{fig:sudakov-ff}
\end{figure}
%%% END fig:sudakov-ff

\cref{fig:sudakov-ff}的图注中 $n$、$\bar n$ 分别是本书的 $n_A$、$n_B$。图画的是类时组态，上例则取类空运动学；两者的区域相同（区域只由各线虚度的标度决定），差别只在最后的解析延拓（\cref{prop:sudakov-oneloop}(b)）。

\begin{remark}[一致性论证：单标度结构、$\mu$ 无关性与极点相消]\label{rem:sudakov-consistency}
\cref{ex:sudakov-offshell}中的双极点系数与双对数可以由三个一般性质确定。(i) \emph{单标度结构}：$H$、$C_i$、$S$ 分别只依赖于 $Q^2$、$P^2$、$P^4/Q^2$
（对超软分母作 $k^\pm\to k^\pm P^2/(2p^\mp)$、$k_T\to k_TP^2/Q$ 的伸缩即见），故 $Q^2R=c_R\,\varepsilon^{-2}(\mu^2/\Lambda_R^2)^\varepsilon(1+O(\varepsilon))$。(ii) 和与 $\mu$ 无关。(iii) 和中极点相消。
记 $a=\ln\frac{\mu^2}{Q^2}$、$b=\ln\frac{\mu^2}{P^2}$，$\ln(\mu^2/\Lambda_R^2)$ 依次为 $a,b,b,2b-a$；$\varepsilon^{-2}$ 相消要求 $c_h+2c_c+c_s=0$，$\varepsilon^{-1}$ 中的对数 $(c_h-c_s)a+2(c_c+c_s)b$
必须为零，故以 $c_h=1$ 归一化有 $(c_h,c_1,c_2,c_s)=(1,-1,-1,1)$，双对数为 $\frac12[a^2-2b^2+(2b-a)^2]=(a-b)^2=\ln^2\frac{Q^2}{P^2}$。
常数 $\pi^2/3$ 不由此决定，它来自显式积分；论证本身又以区域展开成立为前提（\cref{prop:ir-regions-minkowski}）。
\uses{ex:sudakov-offshell}
\end{remark}

\begin{example}[有质量交换：需要解析正规化]\label{ex:sudakov-massive}
令外线在壳且无质量，交换的玻色子有质量 $M$，$t\equiv M^2/Q^2$，并给两条夸克传播子加解析正规子：
\[
I_M=\ee^{\varepsilon\gamma_E}\mu^{2\varepsilon}\!\int\!\frac{\dd^Dk}{\ii\pi^{D/2}}\,\frac{\nu^{2(\delta_1+\delta_2)}}{(-k^2+M^2)\,\big(-(k+p_1)^2\big)^{1+\delta_1}\big(-(k+p_2)^2\big)^{1+\delta_2}} .
\]
$\delta_i=0$、$\varepsilon=0$ 时积分有限，\cite[式~(6.30)]{Jantzen2011}
\begin{equation}\label{eq:sudakov-massive-exact}
Q^2I_M=\frac12\ln^2\frac{Q^2}{M^2}+\ln\frac{Q^2}{M^2}\ln\Big(1-\frac{M^2}{Q^2}\Big)-\mathrm{Li}_2\Big(\frac{M^2}{Q^2}\Big)+\frac{\pi^2}3
\ \xrightarrow{\ \text{领头幂次}\ }\ \frac12\ln^2\frac{Q^2}{M^2}+\frac{\pi^2}3 .
\end{equation}
\emph{纯维数正规化失效。}$\delta_i=0$ 时共线区为 $C_1=\ee^{\varepsilon\gamma_E}\Gamma(1+\varepsilon)\mu^{2\varepsilon}\int_0^1\dd x\,[(1-x)M^2]^{-\varepsilon}/(\varepsilon xQ^2)$，
$x\to0$ 处的 $\int\dd x/x$ 对任何 $\varepsilon$ 都发散：这是\emph{快度发散}（$x$ 是交换玻色子携带的 $p_1^+$ 的分数，横动量固定而 $x\to0$ 时它的快度趋于 $-\infty$），维数正规化只调节横动量积分。
软区（$k\sim M$）在推动 $k^\pm\to\kappa^{\pm1}k^\pm$ 下只差因子 $\kappa^{\delta_1-\delta_2}$：$\delta_1\ne\delta_2$ 时无标度而为零，$\delta_1=\delta_2$ 时无定义。
\emph{加解析正规子。}取 $\delta_1=-\delta_2=\delta$（$\mu=\nu$），每个 $C_i$ 含因子 $\Gamma(\pm2\delta)(Q^2/M^2)^{\pm\delta}$，即 $1/\delta$ 极点（\cite[式~(6.27)]{Jantzen2011}；该文的 $F$ 是这里的 $-I_M$），
它们在和中相消；取 $\mu=Q$，先令 $\delta\to0$ 再展开 $\varepsilon$：
\begin{equation}\label{eq:sudakov-massive-regions}
\lim_{\delta\to0}Q^2(C_1+C_2)=-\frac1{\varepsilon^2}+\frac12\ln^2\frac{Q^2}{M^2}+\frac{5\pi^2}{12}+O(\varepsilon),\qquad
Q^2(H+C_1+C_2)=\frac12\ln^2\frac{Q^2}{M^2}+\frac{\pi^2}3+O(\varepsilon),
% VF-REG-11
\end{equation}
$H$ 即式~\eqref{eq:sudakov-onshell-H}；结果与式~\eqref{eq:sudakov-massive-exact} 的领头幂次相同（\cite[式~(6.29)--(6.30)]{Jantzen2011}，\cite[第~2.3~节，式~(2.50)--(2.52)]{BecherBroggioFerroglia2015}，
后者只用一个正规子：单个区域改变而 $C_1+C_2$ 不变）。正规子须在 $\varepsilon\to0$ 之前移除。
\emph{物理含义。}单个共线积分只知道一个光锥方向与质量 $M$，$C_1+C_2$ 却依赖 $\ln(Q^2/M^2)$：正规子破坏推动不变性，大对数通过 $1/\delta$ 极点的相消重新出现
（$(Q^2/M^2)^{\pm\delta}\Gamma(\pm2\delta)=\pm\frac1{2\delta}+\frac12\ln\frac{Q^2}{M^2}-\gamma_E+O(\delta)$），与 $q_T$ 分布中的“共线反常”\cite{BecherNeubert2011}同源；
快度发散标志着一个新的因子化变量，由演化方程控制（\cref{sec:sudakov-evolution}；\cref{ch:tmd}），正确组合的物理量与正规子无关（参见\cite{BecherBell2012}）。
\emph{对比。}若外线有质量 $p_i^2=m^2$ 而胶子无质量，合并后的质量项为 $x^2m^2$，软端 $\int\dd x\,x^{-1-2\varepsilon}$ 已被维数正规化调节：
$C_1=C_2=-\frac1{Q^2}(\frac{\mu^2}{m^2})^\varepsilon\frac{\ee^{\varepsilon\gamma_E}\Gamma(1+\varepsilon)}{2\varepsilon^2}$，在 $\mu=Q$ 时 $Q^2(H+C_1+C_2)=\frac1\varepsilon\ln\frac{m^2}{Q^2}-\frac12\ln^2\frac{m^2}{Q^2}-\frac{\pi^2}6$，
与精确结果的领头幂次一致。% VF-REG-12
所以只有\emph{交换的玻色子}有质量时才必须使用解析正规子。
\uses{def:ir-regions}
\end{example}

\begin{example}[双对数系数]\label{ex:sudakov-ll}
领头对数只来自既软又共线的区域，此时顶点分子可作程函近似：在近在壳旋量之间 $\gamma^\nu\sla p_2\gamma^\mu\sla p_1\gamma_\nu=-2\sla p_1\gamma^\mu\sla p_2\to4p_1\cdot p_2\,\gamma^\mu$
（旋量恒等式见\cref{prop:sudakov-oneloop}证明的第二步）。于是 $\Lambda^\mu_{\rm LL}=-\ii g^2C_F\cdot4p_1\cdot p_2\,\gamma^\mu\int\frac{\dd^4k}{(2\pi)^4}\frac{1}{k^2(p_1-k)^2(p_2-k)^2}$，
而该积分等于 $-\frac{\ii}{16\pi^2}I$（三个分母各换号，$\dd^4k=\ii\pi^2\cdot\frac{\dd^4k}{\ii\pi^2}$），故 $\delta F_{\rm LL}=-\frac{\alpha_sC_F}{\pi}\,p_1\cdot p_2\,I$，$2p_1\cdot p_2\simeq Q^2$。代入两例的领头幂次：
\begin{equation}\label{eq:sudakov-LL}
\delta F_{\rm LL}^{\text{离壳}}=-\frac{\alpha}{2\pi}\ln\frac{Q^2}{P_1^2}\ln\frac{Q^2}{P_2^2},\qquad
\delta F_{\rm LL}^{M}=-\frac{\alpha}{4\pi}\ln^2\frac{Q^2}{M^2}
% VF-SUD-4
\end{equation}
（前者为离壳外线与无质量光子，后者为在壳外线与质量为 $M$ 的光子），QCD 中 $\alpha\to C_F\alpha_s$。这正是 Collins\cite{Collins:1989bt}所引的 Sudakov 离壳结果 $F\sim\exp[-\frac{g^2}{8\pi^2}\ln^2Q^2]$\cite{Sudakov1956}
与 Jackiw 在壳有质量光子结果 $F\sim\exp[-\frac{g^2}{16\pi^2}\ln^2Q^2]$\cite{Jackiw:1968zz}的一圈系数（$g^2=4\pi\alpha$；参见\cite[第~10.5.2--10.5.3~节]{Collins2011}）。
相差因子 $2$ 的几何原因：在胶子的 $(\ln k^+,\ln k^-)$ 平面中，离壳情形的双对数区域是菱形，有质量光子情形是三角形（$k_T\gtrsim M$ 截去一半），面积比为 $2$（\cite[图~10.7(a)、10.8(b)]{Collins2011}）。
数值示意：$\alpha_sC_F/(2\pi)\approx0.025$（$\alpha_s=0.118$），$\mathcal L=10$ 时离壳单圈双对数已达 $-2.5$，必须重求和；固定耦合下的指数化（\cref{cor:sudakov-exponentiation}）给出抑制因子 $\ee^{-2.5}\approx0.08$。
\uses{ex:sudakov-offshell,ex:sudakov-massive,prop:sudakov-oneloop}
\end{example}
